% Generated draft: requires mathematical and source review.
\documentclass[11pt]{article}
\usepackage[margin=25mm]{geometry}
\usepackage{fontspec}
\setmainfont{Latin Modern Roman}
\usepackage{amsmath,amssymb,mathtools,unicode-math}
\setmathfont{Latin Modern Math}
\usepackage{enumitem,needspace,xurl,hyperref}
\setlength{\parindent}{0pt}
\setlength{\parskip}{5pt}
\setlength{\emergencystretch}{3em}
\newcommand{\N}{\mathbb{N}}
\newcommand{\Z}{\mathbb{Z}}
\newcommand{\Q}{\mathbb{Q}}
\newcommand{\R}{\mathbb{R}}
\newcommand{\C}{\mathbb{C}}
\newcommand{\F}{\mathbb{F}}
\newcommand{\A}{\mathbb{A}}
\newcommand{\PP}{\mathbb{P}}
\newcommand{\E}{\mathbb{E}}
\newcommand{\eps}{\varepsilon}
\newcommand{\dd}{\,\mathrm{d}}
\begin{document}
\section*{UnsolvedMath ID: 30006366}
\textbf{Source problem code:} catalogue entry
\subsection{Definitions and mathematical statement}
The supplied record asks for the multiplicity part of the conjecture called the \emph{\(n\)-or-\(\infty\) conjecture}. Only the lower bound below is supplied; no additional infinite-orbit alternative is asserted here.

For a precise contact-geometric formulation, use the following conventional interpretation of a star-shaped hypersurface. Let \(n\geq 1\) be an integer, and let \(D\subset\R^{2n}\) be a bounded domain containing the origin, with smooth boundary \(M=\partial D\), such that every ray from the origin meets \(M\) exactly once and transversely. In coordinates \((x_1,y_1,\ldots,x_n,y_n)\), put
\[
\lambda_0=\frac12\sum_{j=1}^{n}(x_j\,\dd y_j-y_j\,\dd x_j),\qquad \alpha=\lambda_0|_M.
\]
These hypotheses make \(M\) compact and \(\alpha\) a contact form, meaning \(\alpha\wedge(\dd\alpha)^{n-1}\) is nowhere zero. Its Reeb vector field \(R_\alpha\) is characterized by
\[
\alpha(R_\alpha)=1,\qquad \dd\alpha(R_\alpha,v)=0\quad\text{for every }v\in TM.
\]
A prime closed Reeb orbit is a nonconstant periodic integral curve of \(R_\alpha\), taken with its least positive period, rather than an iterate of a shorter periodic curve. Two such orbits are counted as distinct when their images in \(M\) differ; changes of starting point and multiple traversals do not yield additional orbits.

\textbf{Multiplicity conjecture, under the interpretation above.} For every such \(n\) and \(D\), the Reeb flow of \(\alpha\) has at least \(n\) geometrically distinct prime closed orbits. Equivalently, there exist \(T_1,\ldots,T_n>0\) and periodic integral curves
\[
\gamma_i:\R/T_i\Z\longrightarrow M,\qquad \dot\gamma_i=R_\alpha\circ\gamma_i,
\]
with least periods \(T_i\), whose images are pairwise distinct. No dynamical convexity or nondegeneracy hypothesis is imposed in this conjectural formulation.

\textbf{Unresolved source qualifications.} The source explicitly specifies only a smooth boundary of a star-shaped domain. It does not explicitly specify boundedness, the center of star-shapedness, or transversality of the radial vector field. The formulation above makes the customary compact, strictly star-shaped interpretation explicit; these qualifications require source confirmation, since ordinary set-theoretic star-shapedness alone does not ensure that the displayed restriction is a contact form. The source also emphasizes higher dimensions without giving an explicit range for \(n\); the positive-integer range used above is a stated convention requiring review. Its reference to the standard contact sphere is understood as allowing contact forms defining the standard contact structure, not as fixing a single Reeb flow on the sphere.
\subsection{Short English statement}
The conjecture asks whether every smooth compact strictly star-shaped hypersurface in \(\R^{2n}\), with its induced standard contact form, has at least \(n\) distinct closed Reeb trajectories, counting each trajectory once and excluding repeated traversals. Compactness, strict radial transversality, and the dimension range are explicit interpretive qualifications requiring review. The supplied record labels the general question open and reports a result under dynamical convexity; these are supplied dated assessments, not independently verified conclusions.
\subsection{Sources}
\begin{itemize}
\item Local UnsolvedMath registry: supplied complete record, permanent ID \(30006366\), title \emph{Closed-Orbit Multiplicity on Star-Shaped Hypersurfaces}, exact source problem code \texttt{catalogue entry}. The supplied status is \texttt{open}. Its statement status is \texttt{corrected\_verified}; the accompanying claim that the official report was checked is inherited source metadata, not verification performed for this draft.
\item Supplied report citation: \emph{Dynamische Systeme (2026)}, Oberwolfach Reports, DOI \(10.4171/owr/2025/31\), pp.\ 1647--1706. Supplied source and literature link: \url{https://doi.org/10.4171/owr/2025/31}. The excerpt attributes \emph{Towards the Reeb Hofer--Zehnder and multiplicity conjectures} to Başak Z. Gürel, joint work with Erman Çineli and Viktor L. Ginzburg. It says the report is based on its reference [5], dedicated to the memory of Edi Zehnder, a longtime co-organizer of the workshop; the bibliographic entry for [5] is not supplied.
\item Supplied literature reference, labeled ``Cineli--Ginzburg--Gurel primary paper'': \url{https://arxiv.org/abs/2410.13093}. No full bibliographic citation is supplied, and its identification with the excerpt's reference [5] is not established by the supplied record.
\item Literature assessment dated exactly \texttt{2026-08-22}: the supplied curation labels the general multiplicity conjecture open and says that the recent theorem establishes the lower bound under dynamical convexity, rather than for arbitrary smooth star-shaped hypersurfaces. This is supplied dated curation derived from the OWR record, not independently verified literature assessment; its time-sensitive status and cited claims require independent checking.
\item The historical background names Ekeland, Hofer and Zehnder, Moser, Rabinowitz, Weinstein, and others, and traces the question at least to Lyapunov's time. These are supplied historical attributions without individual bibliographic references.
\end{itemize}
\end{document}
